\documentclass[a4paper]{article}
\usepackage{amsfonts}
\usepackage{amsmath}
\usepackage{amssymb}
\usepackage{graphicx}

\begin{document}

\noindent \large Auteur: Karanjot Singh\\
\noindent \large Studierichting: Informatica\\
\noindent \large Oefeningenreeks 6B nr. 4.4\\

\normalsize

$x_1$ als $x_{15} = 3$ en $v = \dfrac{1}{4}$\\

De formule voor de $n$-de term van een rekenkundige rij is $x_n = x_1 + (n - 1)v$.\\

$x_{15} = x_1 + 14v = 3$\\

$v = \dfrac{1}{4}$\\

$\Rightarrow x_1 + 14 \cdot \dfrac{1}{4} = 3$\\

$\Rightarrow x_1 + \dfrac{14}{4} = 3$\\

$\Rightarrow x_1 + \dfrac{7}{2} = 3$\\

$\Rightarrow x_1 = 3 - \dfrac{7}{2}$\\

$\Rightarrow x_1 = \dfrac{6}{2} - \dfrac{7}{2}$\\

$\Rightarrow x_1 = \dfrac{6 - 7}{2}$\\

$\Rightarrow x_1 = \dfrac{-1}{2}$\\

Dus:\\

$x_1 = \dfrac{-1}{2}$\\

\begin{thebibliography}{999}
	%\bibitem{a} Referentie
\end{thebibliography}

\end{document}
